\documentclass[11pt,a4paper]{article}

\usepackage[utf8]{inputenc}
\usepackage[T1]{fontenc}
\usepackage[a4paper,margin=2.5cm]{geometry}
\usepackage{amsmath, amssymb}
\usepackage{graphicx}
\usepackage{booktabs}
\usepackage{caption}
\usepackage{xcolor}
\usepackage{hyperref}
\usepackage{enumitem}
\usepackage{empheq}

\hypersetup{
    colorlinks=true,
    linkcolor=blue!50!black,
    citecolor=blue!50!black,
    urlcolor=blue!60!black,
    pdftitle={Ultrashort Pulse Propagation Equation -- SI derivation},
}

\definecolor{notecolor}{RGB}{180,120,0}
% Flag for a point that needs checking against the original document.
\newcommand{\checknote}[1]{%
    \par\noindent
    \fcolorbox{notecolor}{notecolor!8}{%
        \parbox{\dimexpr\linewidth-2\fboxsep-2\fboxrule}{%
            \textbf{\textcolor{notecolor}{Check.}} #1%
        }%
    }%
    \par\vspace{0.3em}%
}

\title{Ultrashort Pulse Propagation Equation\\[0.3em]
       \large Full derivation in SI units}
\author{}
\date{}

\begin{document}
\maketitle

\noindent
This document rederives the generalized nonlinear envelope equation entirely in
SI units. It follows the same logical path as the Gaussian-unit derivation, but
every prefactor is recomputed rather than translated, so the two can be compared
independently. Section~\ref{sec:steepening} on self-steepening is written to be
self-contained: it can be edited or lifted out without touching the rest.

\tableofcontents

% =====================================================================
\section{Conventions}
% =====================================================================

\paragraph{Envelope.} The real electric field is written as a carrier at
$\omega_0$ modulated by a slowly varying complex envelope $\mathcal{E}$:
\begin{equation}
\label{eq:envelope-def}
\mathbf{E}(\mathbf{r},t)
 = \tfrac{1}{2}\big[\mathcal{E}(\mathbf{r},t)\,e^{i\psi} + \mathrm{c.c.}\big],
\qquad \psi \equiv k_0 z - \omega_0 t .
\end{equation}
With this convention the peak of the real field is $|\mathcal{E}|$, the cycle
average of its square is $\langle \mathbf{E}^2\rangle = \tfrac12|\mathcal{E}|^2$,
and the intensity is
\begin{equation}
\label{eq:intensity}
\boxed{\;I = \tfrac{1}{2}\,\varepsilon_0 n_0 c\,|\mathcal{E}|^2\;}
\end{equation}
This is the convention used in the solver, where the field array plays the role
of $\mathcal{E}$ and the intensity is formed as \texttt{0.5*eps0*n0*c*abs(E)**2}.
Fixing it once here removes the factor-of-two ambiguities that plague envelope
derivations.

\paragraph{Fourier transform.} We use
\begin{equation}
\mathbf{E}(\mathbf{r},t)=\frac{1}{2\pi}\int \mathbf{E}(\mathbf{r},\omega)\,
   e^{-i\omega t}\,d\omega ,
\end{equation}
so that $\partial_t \to -i\omega$ and, for the detuning $\Omega=\omega-\omega_0$,
\begin{equation}
\label{eq:detuning-rule}
\Omega \;\longleftrightarrow\; i\,\frac{\partial}{\partial t}.
\end{equation}
Every operator below follows from this single rule; getting its sign wrong flips
the sign of both dispersion and self-steepening.

\paragraph{Maxwell (SI).} For a non-magnetic medium ($\mathbf{B}=\mu_0\mathbf{H}$),
\begin{align}
\nabla\cdot\mathbf{D} &= \rho_f, &
\nabla\cdot\mathbf{B} &= 0, \\
\nabla\times\mathbf{E} &= -\frac{\partial \mathbf{B}}{\partial t}, &
\nabla\times\mathbf{H} &= \frac{\partial \mathbf{D}}{\partial t}
   + \mathbf{J}_{\text{ion}} + \mathbf{J}_{\text{pl}} + \mathbf{J}_{\text{STE}},
\end{align}
with $\mathbf{D}=\varepsilon_0\mathbf{E}+\mathbf{P}$ and
$\mu_0\varepsilon_0 = 1/c^2$.

\checknote{In the Gaussian document, Amp\`ere's law is written
$\nabla\times\tilde{\mathbf H} = \nabla\times\tilde{\mathbf B}$. That equality
holds only because $\mu=1$ in Gaussian units; it is worth stating explicitly,
since it is the step that assumes a non-magnetic medium. Also, $\mathbf{J}_{\text{STE}}$
appears in the wave equation there but is absent from Amp\`ere's law -- it should
be introduced in both places.}

% =====================================================================
\section{From Maxwell to the wave equation}
% =====================================================================

Taking the curl of Faraday's law and using
$\nabla\times(\nabla\times\mathbf{A})=\nabla(\nabla\cdot\mathbf{A})-\nabla^2\mathbf{A}$
with $\nabla\cdot\mathbf{E}\simeq 0$ (homogeneous medium):
\begin{equation*}
-\nabla^2\mathbf{E}
 = -\mu_0\frac{\partial}{\partial t}
   \Big[\frac{\partial\mathbf{D}}{\partial t}+\mathbf{J}\Big],
\qquad \mathbf{J}\equiv\mathbf{J}_{\text{ion}}+\mathbf{J}_{\text{pl}}+\mathbf{J}_{\text{STE}} .
\end{equation*}
Splitting the polarization into linear and nonlinear parts,
$\mathbf{D}=\mathbf{D}^{(1)}+\mathbf{P}_{\text{NL}}$ with
$\mathbf{D}^{(1)}=\varepsilon_0\mathbf{E}+\mathbf{P}^{(1)}$:
\begin{equation}
\label{eq:wave-time}
\nabla^2\mathbf{E}-\mu_0\frac{\partial^2\mathbf{D}^{(1)}}{\partial t^2}
 = \mu_0\frac{\partial^2\mathbf{P}_{\text{NL}}}{\partial t^2}
 + \mu_0\frac{\partial\mathbf{J}}{\partial t}.
\end{equation}
Comparing with the Gaussian form, the entire unit conversion is the single
substitution
\begin{equation}
\label{eq:conversion}
\frac{4\pi}{c^2}\;\longrightarrow\;\mu_0=\frac{1}{\varepsilon_0 c^2}
\qquad\text{i.e.}\qquad 4\pi \longrightarrow \frac{1}{\varepsilon_0}.
\end{equation}

In the frequency domain, with
$\mathbf{D}^{(1)}(\omega)=\varepsilon_0\varepsilon_r(\omega)\mathbf{E}(\omega)$,
\begin{equation}
\nabla^2\mathbf{E}+\frac{\omega^2}{c^2}\varepsilon_r(\omega)\mathbf{E}
 = -\mu_0\omega^2\underbrace{\Big[\mathbf{P}_{\text{NL}}
    + \frac{i\mathbf{J}}{\omega}\Big]}_{\displaystyle \equiv\;\mathcal{S}(\mathbf{r},\omega)} .
\end{equation}
The grouping $\mathcal{S}$ plays the role of the symbol $\text{\it ka}$ in the
Gaussian document.

% =====================================================================
\section{Envelope and dispersion}
% =====================================================================

Inserting $\mathbf{E}(\mathbf{r},\omega)\simeq
\mathcal{E}(\mathbf{r},\Omega)e^{ik_0z}$ and expanding the Laplacian,
\begin{equation*}
\Big[\nabla_\perp^2+\partial_z^2+2ik_0\partial_z-k_0^2\Big]\mathcal{E}
 + k^2(\omega)\,\mathcal{E} = -\mu_0\omega^2\,\mathcal{S}\,e^{-ik_0z},
\end{equation*}
with $k^2(\omega)=\varepsilon_r(\omega)\omega^2/c^2$. Expanding $k(\omega)$
around $\omega_0$,
\begin{equation}
k(\omega)=k_0+k_1\Omega+\underbrace{\sum_{n\ge 2}\frac{k_n}{n!}\Omega^n}_{\displaystyle \equiv\, L},
\qquad k_1=\frac{1}{v_g}=\frac{n_g}{c},
\end{equation}
so that $k^2-k_0^2 = 2k_0k_1\Omega + k_1^2\Omega^2 + 2k_0L + 2k_1 L\Omega + L^2$.
Applying rule~\eqref{eq:detuning-rule} and moving to the retarded frame
$\tau=t-k_1z$, $z'=z$ (which removes the $2k_0k_1\Omega$ term), and dropping
$\partial_{z'}^2$ (SVEA):
\begin{equation}
\label{eq:pre-U}
\Big[\nabla_\perp^2
 + 2ik_0\partial_{z'}\big(1+\tfrac{ik_1}{k_0}\partial_\tau\big)
 + 2k_0\hat{L}\big(1+\tfrac{ik_1}{k_0}\partial_\tau\big)
 + \hat{L}^2\Big]\mathcal{E}
 = -\mu_0\omega_0^2\,\hat{T}^2\mathcal{P}_{\text{NL}} + \dots
\end{equation}
where $\hat L=\sum_{n\ge2}\frac{k_n}{n!}(i\partial_\tau)^n$.

% =====================================================================
\section{The two operators \texorpdfstring{$\hat U$}{U} and \texorpdfstring{$\hat T$}{T}}
\label{sec:operators}
% =====================================================================

Two distinct first-order temporal operators appear, and they are easy to
conflate because they look alike.

\paragraph{Space--time focusing, $\hat U$.} The bracket generated by the
retarded-frame change of variables in~\eqref{eq:pre-U} is
\begin{equation}
\label{eq:U-def}
\hat U \;=\; 1+\frac{k_1}{k_0}\,\Omega
 \;=\; 1+\frac{n_g}{n_0}\frac{\Omega}{\omega_0}
 \;=\; 1+i\,\frac{n_g}{n_0\omega_0}\frac{\partial}{\partial\tau},
\end{equation}
using $k_1/k_0=(n_g/c)(c/n_0\omega_0)=n_g/(n_0\omega_0)$. It multiplies
$\partial_{z'}$ and the dispersion operator; physically it says that different
frequency components diffract differently, which couples the transverse and
temporal dynamics.

\paragraph{Self-steepening, $\hat T$.} This one has a completely different
origin: it is what the \emph{second time derivative} of the nonlinear
polarization produces. With
$\mathbf{P}_{\text{NL}}=\tfrac12\mathcal{P}_{\text{NL}}e^{i\psi}+\mathrm{c.c.}$,
\begin{align}
\frac{\partial^2}{\partial t^2}\Big[\mathcal{P}_{\text{NL}}e^{-i\omega_0 t}\Big]
&= \Big[\ddot{\mathcal{P}}-2i\omega_0\dot{\mathcal{P}}-\omega_0^2\mathcal{P}\Big]e^{-i\omega_0t}
\nonumber\\
&= -\omega_0^2\Big[1+\frac{2i}{\omega_0}\partial_t-\frac{1}{\omega_0^2}\partial_t^2\Big]
   \mathcal{P}\,e^{-i\omega_0t}
 = -\omega_0^2\,\hat T^2\,\mathcal{P}\,e^{-i\omega_0 t},
\end{align}
so the square is exact, not an approximation:
\begin{equation}
\label{eq:T-def}
\hat T = 1+\frac{i}{\omega_0}\frac{\partial}{\partial\tau}.
\end{equation}

\paragraph{Why the literature writes a single $\hat T$.} Couairon and
Mysyrowicz write the Kerr term with one power of $\hat T$, not two. The two
statements are consistent. Dividing the propagation equation through by $\hat U$
leaves the Kerr term proportional to $\hat T^2/\hat U$, and comparing
\eqref{eq:U-def} with \eqref{eq:T-def},
\begin{equation}
\hat U = 1+i\frac{n_g}{n_0}\frac{1}{\omega_0}\partial_\tau,
\qquad
\hat T = 1+i\frac{1}{\omega_0}\partial_\tau
\qquad\Longrightarrow\qquad
\hat U \simeq \hat T \;\;\text{when}\;\; n_g\simeq n_0 .
\end{equation}
For fused silica at $800$~nm, $n_0=1.4533$ and $n_g\simeq 1.4671$, so
$n_g/n_0=1.0095$: the two operators agree to within $1\%$. Hence
\begin{equation}
\label{eq:T2overU}
\frac{\hat T^2}{\hat U}\;\simeq\;\hat T ,
\end{equation}
and the single-$\hat T$ form of the literature is recovered. This is the
cleanest way to reconcile the two conventions, and it matters for
Section~\ref{sec:steepening}: the physically meaningful shock coefficient is the
one obtained \emph{after} dividing by $\hat U$.

% =====================================================================
\section{Nonlinear polarization and \texorpdfstring{$n_2$}{n2}}
% =====================================================================

In SI, $\mathbf{P}=\varepsilon_0\chi^{(1)}\mathbf{E}+\varepsilon_0\chi^{(3)}\mathbf{E}^3$.
Cubing~\eqref{eq:envelope-def},
\begin{equation}
\mathbf{E}^3=\tfrac18\mathcal{E}^3e^{3i\psi}
 +\tfrac38|\mathcal{E}|^2\mathcal{E}\,e^{i\psi}+\mathrm{c.c.}
\end{equation}
The $3\omega_0$ term drives third-harmonic generation and is dropped. Writing
$\mathbf{P}_{\text{NL}}=\tfrac12\mathcal{P}_{\text{NL}}e^{i\psi}+\mathrm{c.c.}$,
\begin{equation}
\label{eq:PNL}
\mathcal{P}_{\text{NL}}=\tfrac34\,\varepsilon_0\chi^{(3)}|\mathcal{E}|^2\mathcal{E}.
\end{equation}
The refractive index follows from
$\varepsilon_r=n_0^2+\tfrac34\chi^{(3)}|\mathcal{E}|^2$:
\begin{equation}
n=\sqrt{\varepsilon_r}\simeq n_0+\frac{3\chi^{(3)}}{8n_0}|\mathcal{E}|^2
 \equiv n_0+n_2 I ,
\end{equation}
and inserting $I$ from~\eqref{eq:intensity},
\begin{equation}
\label{eq:n2}
\boxed{\;n_2=\frac{3\chi^{(3)}}{4\,\varepsilon_0 n_0^2 c}
\qquad\Longleftrightarrow\qquad
\chi^{(3)}=\tfrac43\,\varepsilon_0 c\, n_0^2\, n_2 \;}
\end{equation}
which is the relation already used in the main report. Note that
$\chi^{(3)}$ is \emph{not} numerically the same quantity as in Gaussian units:
the two differ both by $\varepsilon_0$ (from $\mathbf{P}=\varepsilon_0\chi^{(3)}\mathbf{E}^3$
versus $\tilde{\mathbf{P}}=\chi^{(3)}\tilde{\mathbf{E}}^3$) and by the $4\pi$ in
$\mathbf{D}$. Only $n_2$, being defined through the measurable
$\Delta n = n_2 I$, is convention-independent.

\paragraph{Kerr source term.} Substituting~\eqref{eq:PNL} and~\eqref{eq:n2} into
the right-hand side of~\eqref{eq:pre-U}, then dividing by $2ik_0$:
\begin{align}
i\frac{\mu_0\omega_0^2}{2k_0}\hat T^2\mathcal{P}_{\text{NL}}
&= i\frac{\mu_0\omega_0^2}{2k_0}\,\varepsilon_0^2 n_0^2 c\, n_2\,
   \hat T^2\big[|\mathcal{E}|^2\mathcal{E}\big]
\nonumber\\
&= i\frac{\omega_0^2\varepsilon_0 n_0^2 n_2}{2k_0 c}\,
   \hat T^2\big[|\mathcal{E}|^2\mathcal{E}\big]
   \qquad (\mu_0\varepsilon_0^2=\varepsilon_0/c^2)
\nonumber\\
&= i\,\frac{\omega_0}{c}\,n_2\,\hat T^2\big[I\,\mathcal{E}\big]
   \qquad (k_0=n_0\omega_0/c,\;\; \varepsilon_0 n_0|\mathcal{E}|^2=2I/c),
\end{align}
that is,
\begin{empheq}[box=\fbox]{equation}
\label{eq:kerr-term}
\text{Kerr}\;=\;i\,K_0\,n_2\,\hat T^2\big[I\,\mathcal{E}\big],
\qquad K_0\equiv\frac{\omega_0}{c}\;\;(\text{vacuum wavenumber}).
\end{empheq}
This is \emph{identical} to the Gaussian result. The agreement is a genuine
cross-check: the $4\pi$'s and the $\varepsilon_0$'s cancel because $n_2$ is
defined through the intensity in both systems.

% =====================================================================
\section{Self-steepening}
\label{sec:steepening}
% =====================================================================

\noindent
\emph{This section is self-contained. It replaces the qualitative paragraph
``the more intense peak travels more slowly than the wings'' by an explicit
derivation, and identifies exactly which part of that statement is right and
which part is hand-waving.}

\subsection{Extracting the shock term}

Start from the Kerr term~\eqref{eq:kerr-term} after division by $\hat U$, which
by~\eqref{eq:T2overU} leaves a single $\hat T$. Keep only the Kerr nonlinearity
(no diffraction, no dispersion, no plasma) so that the effect is isolated:
\begin{equation}
\label{eq:shock-model}
\frac{\partial\mathcal{E}}{\partial z}
 = i\,K_0 n_2\,\hat T\big[I\mathcal{E}\big]
 = \underbrace{i\,K_0 n_2\, I\,\mathcal{E}}_{\text{SPM}}
   \;\underbrace{-\;\frac{n_2}{c}\,\frac{\partial}{\partial\tau}\big(I\mathcal{E}\big)}_{\text{shock}},
\end{equation}
where the second term used $\hat T = 1+\tfrac{i}{\omega_0}\partial_\tau$ and
$K_0/\omega_0=1/c$. Define the shock coefficient
\begin{equation}
s \equiv \frac{n_2}{c} .
\end{equation}
(If one keeps $\hat T^2$ without dividing by $\hat U$, the coefficient is $2s$;
this is the origin of the $2n_2/c$ quoted in the report. The physically
meaningful value, after $\hat U$ has been divided out, is $s$.)

\subsection{The intensity obeys a Burgers equation}

Write the envelope in amplitude--phase form, $\mathcal{E}=\sqrt{I}\,e^{i\varphi}$
(absorbing the constant of~\eqref{eq:intensity}, which does not affect the
structure). Then
\begin{equation*}
\frac{\partial\mathcal{E}}{\partial z}
 = \Big[\partial_z\sqrt{I}+i\sqrt{I}\,\partial_z\varphi\Big]e^{i\varphi},
\qquad
\frac{\partial}{\partial\tau}\big(I\mathcal{E}\big)
 = \Big[\partial_\tau\big(I^{3/2}\big)+iI^{3/2}\partial_\tau\varphi\Big]e^{i\varphi}.
\end{equation*}
Separating real and imaginary parts of~\eqref{eq:shock-model}:
\begin{align}
\text{(real)}\quad
 &\partial_z\sqrt{I} = -s\,\partial_\tau\big(I^{3/2}\big)
   = -\tfrac32 s\sqrt{I}\,\partial_\tau I , \\
\text{(imag)}\quad
 &\partial_z\varphi = K_0n_2 I - s\,I\,\partial_\tau\varphi .
\end{align}
Multiplying the real equation by $2\sqrt{I}$ gives the central result:
\begin{empheq}[box=\fbox]{equation}
\label{eq:burgers}
\frac{\partial I}{\partial z} + 3s\,I\,\frac{\partial I}{\partial\tau} = 0,
\qquad s=\frac{n_2}{c}.
\end{empheq}

This is the \textbf{inviscid Burgers equation}, the archetype of shock
formation. It is what justifies the word ``steepening'', and everything
qualitative about self-steepening follows from it without further argument:

\begin{enumerate}[leftmargin=1.6em]
\item[(i)] \textbf{Each intensity level travels at its own speed.} Along the
characteristics $d\tau/dz = 3sI$, the intensity is constant. High-intensity
points drift toward \emph{larger} $\tau$. Since $\tau=t-z/v_g$, larger $\tau$
means arriving later: the peak is retarded relative to the wings. The
statement ``the peak travels more slowly'' is therefore correct.

\item[(ii)] \textbf{The distortion is intrinsically asymmetric.} The advection
velocity depends on $I$, not on $\partial_\tau I$. The peak overtakes the
trailing edge (which lies at larger $\tau$) and runs away from the leading edge.
The trailing edge compresses and steepens; the leading edge flattens. A
perfectly symmetric Gaussian input therefore becomes asymmetric immediately,
which is the point the qualitative paragraph was trying to make.

\item[(iii)] \textbf{A shock forms at finite distance} (next subsection).
\end{enumerate}

\subsection{Shock distance}

For $\partial_z I + c(I)\partial_\tau I = 0$ with $c(I)=3sI$, characteristics
cross, and $\partial_\tau I$ becomes infinite, at
\begin{equation}
z_{\text{sh}}
 = \frac{-1}{\displaystyle\min_\tau \frac{\partial c(I_0)}{\partial\tau}}
 = \frac{-1}{3s\,\displaystyle\min_\tau \partial_\tau I_0}.
\end{equation}
For a Gaussian input $I_0(\tau)=I_p\exp(-\tau^2/t_p^2)$, the steepest negative
slope sits at $\tau=t_p/\sqrt2$ and equals
$\min_\tau\partial_\tau I_0=-\sqrt{2}\,e^{-1/2}\,I_p/t_p\simeq-0.858\,I_p/t_p$.
Hence
\begin{empheq}[box=\fbox]{equation}
\label{eq:zshock}
z_{\text{sh}} \;=\; \frac{1}{3\sqrt2 e^{-1/2}}\,\frac{c\,t_p}{n_2 I_p}
 \;\simeq\; 0.39\,\frac{c\,t_p}{n_2 I_p}.
\end{empheq}
The numerical coefficient $0.39$ reproduces the classical result of Anderson and
Lisak, quoted by Agrawal as $z_{\text{sh}}=0.39\,L_{\text{NL}}/s'$ with
$L_{\text{NL}}=1/(K_0n_2I_p)$ and $s'=1/(\omega_0 t_p)$. Recovering it from the
Burgers picture is a useful check that the factor $3$
in~\eqref{eq:burgers} is right.

\paragraph{How far is that, in practice?} Comparing to the self-focusing length
$L_{\text{SF}}=1/(n_2k_0I_p)$ of the same beam,
\begin{equation}
\frac{z_{\text{sh}}}{L_{\text{SF}}} = 0.39\,n_0\,\omega_0 t_p .
\end{equation}
For fused silica at $800$~nm with $t_p=136$~fs, $\omega_0t_p\simeq320$ and
$n_0=1.4533$, giving $z_{\text{sh}}\simeq 180\,L_{\text{SF}}$. Self-steepening
is thus a very slow effect compared with self-focusing: on its own it would
never matter. It becomes visible in filamentation only because the collapse
first compresses the pulse in time (small $t_p$) and raises the intensity (large
$I_p$), both of which shorten $z_{\text{sh}}$ dramatically. This is why the
shock term is negligible on the leading edge of the propagation and abruptly
important near the nonlinear focus.

\subsection{Spectral signature: why the blue side runs away}

From the imaginary part, to leading order the phase is the ordinary SPM phase
$\varphi \simeq K_0 n_2 I z$, so the instantaneous frequency shift is
\begin{equation}
\label{eq:freqshift}
\delta\omega(\tau) = -\frac{\partial\varphi}{\partial\tau}
 = -K_0 n_2 z\,\frac{\partial I}{\partial\tau}.
\end{equation}
On the leading edge $\partial_\tau I>0$, so $\delta\omega<0$: red shift. On the
trailing edge $\partial_\tau I<0$, so $\delta\omega>0$: blue shift. Without the
shock term the two edges are mirror images and the spectrum is symmetric.

With the shock term, \eqref{eq:burgers} makes $|\partial_\tau I|$ on the
trailing edge grow without bound as $z\to z_{\text{sh}}$, while the leading edge
flattens. Through~\eqref{eq:freqshift} the blue edge of the spectrum therefore
extends much further than the red edge, and diverges at the shock. This is the
mechanism behind the blue-shifted asymmetry of supercontinuum
spectra~\cite{rothenberg1992,gaeta2000}, and it is a \emph{prediction} of
\eqref{eq:burgers} rather than an assertion.

\subsection{What the ``intensity-dependent group velocity'' argument gets wrong}

The usual one-line justification is: the nonlinear index adds $n_2I$ to the
group index, so $1/v_g \to (n_g+n_2I)/c$, and in the retarded frame a point of
intensity $I$ drifts at
\begin{equation}
\left.\frac{d\tau}{dz}\right|_{\text{naive}} = \frac{n_2 I}{c} = s\,I .
\end{equation}
Compare with the exact characteristic velocity of~\eqref{eq:burgers},
$d\tau/dz = 3sI$. The naive argument gets the \emph{sign} right (peak retarded)
and the \emph{scaling} right (linear in $I$), but underestimates the drift by a
factor $3$. The missing factor comes from the shock operator acting on the
product $I\mathcal{E}$, not on the phase alone: differentiating $I^{3/2}$
produces $\tfrac32 I^{1/2}\partial_\tau I$, and converting back from
$\sqrt I$ to $I$ contributes the remaining factor $2$.

So the qualitative picture in the report is not wrong, but it should not be
presented as a derivation. The honest formulation is: \emph{the shock term makes
the intensity obey a Burgers equation whose characteristic velocity is
$3n_2I/c$; an intensity-dependent group velocity is a useful mental image for
the sign of the effect, but it does not give the correct coefficient.}

\checknote{The report currently writes the shock term as
$(2n_2/c)\,\partial(|\mathcal{E}|^2\mathcal{E})/\partial t$. The factor $2$ is
consistent with keeping $\hat T^2$ and using $K_0$ rather than $k_0$, but the
quantity being differentiated should be the \emph{intensity} times the envelope,
$I\mathcal{E}$, not $|\mathcal{E}|^2\mathcal{E}$ -- the two differ by
$\tfrac12\varepsilon_0n_0c$. Worth fixing when the paragraph is rewritten.}

\subsection{Numerical verification}

Everything above is checkable without the full solver, since
\eqref{eq:shock-model} is a one-dimensional equation. The module
\texttt{sim/self\_steepening.py} integrates it with RK4 and a spectral
$\partial_\tau$, and produces two figures.

\paragraph{The Burgers law, tested rather than asserted.} Differentiating
\eqref{eq:burgers} along a characteristic gives $dJ/dz=-3sJ^2$ for
$J=\partial_\tau I$, hence
\begin{equation}
\label{eq:slopelaw}
\frac{|J|(z)}{|J|(0)} = \frac{1}{1-z/z_{\text{sh}}} .
\end{equation}
A threshold at $G$ times the initial slope must therefore be crossed at exactly
$z=(1-1/G)\,z_{\text{sh}}$, independently of $t_p$, $I_p$ and $n_2$. Measuring
that crossing for $G=3$ over input durations $t_p=15$ to $80$~fs gives
$z_{\text{meas}}/z_{\text{sh}}=0.6667$ with a maximum deviation of
$1\times10^{-4}$ from the predicted $2/3$, and the slope-growth curves for all
six durations collapse onto \eqref{eq:slopelaw}. This validates the factor $3$
in \eqref{eq:burgers} and the coefficient $0.39$ in \eqref{eq:zshock}
simultaneously.

\paragraph{The spectral asymmetry.} Propagating a $30$~fs, $5\times10^{13}$
W/cm$^2$ pulse to $0.85\,z_{\text{sh}}$:
\begin{center}
\begin{tabular}{lcc}
\toprule
 & shock OFF & shock ON \\
\midrule
trailing-edge slope $\partial_\tau I$ (norm.) & $-0.0286$ & $-0.1901$ \\
leading-edge slope $\partial_\tau I$ (norm.)  & $+0.0286$ & $+0.0155$ \\
instantaneous frequency range (rad/fs) & $[-0.67,+0.67]$ & $[-0.43,+1.62]$ \\
spectral extent, blue/red at $10^{-3}$ & $1.00$ & $3.68$ \\
\bottomrule
\end{tabular}
\end{center}
Without the shock term the pulse shape is rigorously unchanged (SPM is a pure
phase) and every quantity is symmetric to machine precision. With it, the
trailing edge steepens by $6.6\times$ while the leading edge flattens, exactly
as \eqref{eq:burgers} requires, and the blue edge of the spectrum runs
$3.7\times$ further than the red one.

\paragraph{A caution on the metric.} The natural-looking measure -- the ratio of
spectral \emph{energy} above and below $\omega_0$ -- is the wrong observable
here: it comes out slightly \emph{red}-heavy ($0.75$) even in the case where the
blue edge extends four times further. The reason is that the blue tail is
generated by a single narrow region of $\tau$ (the steepened edge), so it is
spread thinly over a wide bandwidth, whereas the bulk of the energy stays in the
central SPM peaks. The asymmetry claimed here is one of spectral \emph{extent},
not of energy, and the figure reports the extent.

% =====================================================================
\section{Delayed Raman response}
% =====================================================================

The Raman derivation is unit-independent once the driving term is normalized, so
it carries over unchanged. A Raman-active vibrational mode of coordinate $q$,
frequency $\Omega_R$ and damping time $\tau_d$ is driven by the cycle-averaged
field through the Placzek polarizability expansion
$\alpha(q)\simeq\alpha_0+\alpha' q$:
\begin{equation}
\mu\Big[\frac{d^2}{dt^2}+\frac{2}{\tau_d}\frac{d}{dt}+\Omega_R^2\Big]q
 = \tfrac12\alpha'\langle \mathbf{E}^2\rangle
 = \tfrac14\alpha'|\mathcal{E}|^2 .
\end{equation}
Normalizing to $Q\equiv q\,\Omega_R^2/\kappa$ with $\kappa=\alpha'/(4\mu)$, and
expressing the drive as an intensity,
\begin{equation}
\Big[\frac{d^2}{dt^2}+\frac{2}{\tau_d}\frac{d}{dt}+\Omega_R^2\Big]Q
 = \Omega_R^2\, I .
\end{equation}
The Green function of this operator, causal and normalized by the jump
condition $G'(0^+)-G'(0^-)=1$, is
\begin{equation}
G(t)=\tau_s\,e^{-t/\tau_d}\sin\!\big(t/\tau_s\big)\,\mathcal{H}(t),
\qquad
\frac{1}{\tau_s}=\sqrt{\Omega_R^2-\frac{1}{\tau_d^2}} ,
\end{equation}
so the response function is
\begin{equation}
\label{eq:raman-R}
R(t)=\Omega_R^2\,\tau_s\,e^{-t/\tau_d}\sin\!\big(t/\tau_s\big)
 = \frac{\tau_s^2+\tau_d^2}{\tau_s\tau_d^2}\,e^{-t/\tau_d}\sin\!\big(t/\tau_s\big),
\end{equation}
using $\Omega_R^2=\tau_s^{-2}+\tau_d^{-2}$. The second form is the one usually
quoted in the filamentation literature. It is correctly normalized:
\begin{equation}
\int_0^\infty R(t)\,dt
 = \Omega_R^2\tau_s\cdot\frac{1/\tau_s}{\tau_d^{-2}+\tau_s^{-2}} = 1 .
\end{equation}
The Kerr response is then split between an instantaneous and a delayed part with
weight $f_R$:
\begin{equation}
\label{eq:kerr-full}
I \;\longrightarrow\;
 (1-f_R)\,I(\mathbf{r},t) + f_R\!\int_{-\infty}^{t}\! R(t-\tau)\,I(\mathbf{r},\tau)\,d\tau .
\end{equation}
The solver uses $f_R=0.18$, $\tau_d=32$~fs, $\tau_s=12$~fs.

\checknote{In the Gaussian document the normalized equation is written once with
$\Omega^2|\mathcal{E}|^2$ on the right and once with $\Omega^2 I$; these differ
by $n_0c/8\pi$. Only the second is consistent with $R$ being normalized to unit
area and with $I$ entering~\eqref{eq:kerr-full}. The $Q$ substitution is also
performed twice in a row.}

% =====================================================================
\section{Plasma response: the Drude model}
% =====================================================================

A conduction electron of effective mass $m_{\text{eff}}$, driven by the field and
damped by collisions at rate $1/\tau_c$:
\begin{equation}
m_{\text{eff}}\frac{d\mathbf{v}}{dt}
 = -e\mathbf{E}-\frac{m_{\text{eff}}}{\tau_c}\mathbf{v}.
\end{equation}
Seeking $\mathbf{v}=\tfrac12\tilde{\mathbf{v}}e^{i\psi}+\mathrm{c.c.}$ and using
the SVEA ($|\partial_t\tilde{\mathbf v}|\ll\omega_0|\tilde{\mathbf v}|$, so
$d\mathbf{v}/dt\simeq-i\omega_0\tilde{\mathbf v}$):
\begin{equation}
\tilde{\mathbf v}
 = -\frac{e\tau_c}{m_{\text{eff}}(1-i\omega_0\tau_c)}\,\mathcal{E}
\quad\Longrightarrow\quad
\mathcal{J}_{\text{pl}} = -e\rho\,\tilde{\mathbf v}
 = \frac{e^2\rho\,\tau_c\,(1+i\omega_0\tau_c)}{m_{\text{eff}}(1+\omega_0^2\tau_c^2)}\,\mathcal{E}.
\end{equation}
The corresponding source term in the propagation equation is
\begin{align}
\frac{\mu_0}{2ik_0}\frac{\partial \mathcal{J}_{\text{pl}}}{\partial t}
&\simeq \frac{\mu_0}{2ik_0}(-i\omega_0)\mathcal{J}_{\text{pl}}
 = -\frac{\mu_0\omega_0}{2k_0}\frac{e^2\tau_c}{m_{\text{eff}}(1+\omega_0^2\tau_c^2)}
   \,\rho\,(1+i\omega_0\tau_c)\,\mathcal{E}
\nonumber\\
&= -\frac{e^2\omega_0\tau_c}
        {2k_0\varepsilon_0c^2 m_{\text{eff}}(1+\omega_0^2\tau_c^2)}
   \,\rho\,(1+i\omega_0\tau_c)\,\mathcal{E}.
\end{align}
Defining the inverse-bremsstrahlung cross section in SI,
\begin{empheq}[box=\fbox]{equation}
\label{eq:sigma}
\sigma = \frac{k_0\,e^2}{n_0^2\,\omega_0^2\,\varepsilon_0\,m_{\text{eff}}}
         \cdot\frac{\omega_0\tau_c}{1+\omega_0^2\tau_c^2},
\end{empheq}
and using $n_0^2\omega_0^2=k_0^2c^2$, the prefactor above is exactly $\sigma/2$:
\begin{equation}
\frac{\sigma}{2}
 = \frac{e^2\,\omega_0\tau_c}{2k_0c^2\varepsilon_0 m_{\text{eff}}(1+\omega_0^2\tau_c^2)} .
\end{equation}
Hence
\begin{equation}
\label{eq:plasma-term}
\text{plasma}\;=\;-\frac{\sigma}{2}\,(1+i\omega_0\tau_c)\,\rho\,\mathcal{E},
\end{equation}
the real part being absorption (inverse bremsstrahlung) and the imaginary part
the defocusing phase. Expression~\eqref{eq:sigma} is exactly what the solver
computes:
\begin{center}
\texttt{sigmaomega = k0*q\_e**2*tau\_c / (n0**2 * m * eps0 * omega0 * (1+(omega0*tau\_c)**2))}
\end{center}
so this equation is the direct justification of that line of code. No Gaussian
detour is needed.

% =====================================================================
\section{Photoionization losses}
% =====================================================================

Energy conservation (Poynting) requires that the power transferred to the newly
freed electrons equals the gap energy times the ionization rate:
\begin{equation}
\frac{\partial u}{\partial t}+\nabla\!\cdot\!\mathbf{S}
 = -\mathbf{J}_{\text{ion}}\!\cdot\!\mathbf{E},
\qquad
\big\langle \mathbf{J}_{\text{ion}}\!\cdot\!\mathbf{E}\big\rangle_T
 = U_i\,W_{\text{PI}} .
\end{equation}
Taking the current in phase with the field, $\mathbf{J}_{\text{ion}}=\gamma_{\text{ion}}\mathbf{E}$,
and using $\langle\mathbf{E}^2\rangle=\tfrac12|\mathcal{E}|^2$:
\begin{equation}
\gamma_{\text{ion}}\,\tfrac12|\mathcal{E}|^2 = U_iW_{\text{PI}}
\quad\Longrightarrow\quad
\gamma_{\text{ion}} = \frac{2U_iW_{\text{PI}}}{|\mathcal{E}|^2}.
\end{equation}
The source term follows the same route as the plasma one, keeping a single
$\hat T$ because only one time derivative acts on $\mathbf{J}$:
\begin{align}
-\frac{\mu_0\omega_0}{2k_0}\hat T\big[\gamma_{\text{ion}}\mathcal{E}\big]
&= -\frac{\mu_0\omega_0}{2k_0}\,\hat T\!\left[\frac{2U_iW_{\text{PI}}}{|\mathcal{E}|^2}\mathcal{E}\right]
\nonumber\\
&= -\frac{\mu_0\omega_0}{2k_0}\,\varepsilon_0 n_0 c\;
    \hat T\!\left[\frac{U_iW_{\text{PI}}}{I}\mathcal{E}\right]
&&\big(|\mathcal{E}|^2 = 2I/\varepsilon_0n_0c\big)
\nonumber\\
&= -\frac{1}{2}\,\hat T\!\left[\frac{U_iW_{\text{PI}}}{I}\mathcal{E}\right]
&&\big(\mu_0\varepsilon_0 n_0 c\,\omega_0/k_0 = 1\big).
\end{align}
The factor $\tfrac12$ therefore comes out identical to the Gaussian derivation,
confirming that it is not a units artifact. Dimensionally
$[U_iW_{\text{PI}}/I]=\text{m}^{-1}$, as an absorption coefficient must be: the
denominator is an \emph{intensity}, not a squared field amplitude.

The rate $W_{\text{PI}}$ is the Keldysh photoionization rate,
\begin{equation}
W_{\text{PI}}
 = \frac{2\omega_0}{9\pi}\left(\frac{\omega_0 m}{\hbar\sqrt{\Gamma_1}}\right)^{3/2}
   \left[\sqrt{\frac{\pi}{2K(\Gamma_2)}}\sum_{n=0}^{\infty}e^{-n\alpha}\,
   \Phi\!\left(\sqrt{\beta(n+2\nu)}\right)\right] e^{-\alpha\langle x+1\rangle},
\end{equation}
with $\gamma=\omega_0\sqrt{mU_i}/(e|\mathcal{E}|)$,
$\Gamma_1=\gamma^2/(1+\gamma^2)$, $\Gamma_2=1/(1+\gamma^2)$,
$\alpha=\pi[K(\Gamma_1)-E(\Gamma_1)]/E(\Gamma_2)$,
$\beta=\pi^2/[4K(\Gamma_2)E(\Gamma_2)]$,
$x=(2U_i/\pi\hbar\omega_0)\,E(\Gamma_2)/\sqrt{\Gamma_1}$ and
$\nu=\langle x+1\rangle-x$; $K$ and $E$ are the complete elliptic integrals
(argument $m=k^2$) and $\Phi$ the Dawson function. The limits $\gamma\gg1$ and
$\gamma\ll1$ give multiphoton and tunnel ionization respectively.

% =====================================================================
\section{The complete equation}
% =====================================================================

Collecting~\eqref{eq:kerr-term}, \eqref{eq:kerr-full}, \eqref{eq:plasma-term}
and the ionization term:
\begin{empheq}[box=\fbox]{equation}
\label{eq:master}
\begin{split}
\hat U\,\frac{\partial \mathcal{E}}{\partial z'}
 =\;& i\left[\frac{\nabla_\perp^2}{2k_0}
     +\hat L\Big(\hat U+\frac{\hat L}{2k_0}\Big)\right]\mathcal{E} \\[0.3em]
 &+ i\,K_0 n_2\,\hat T^2\!\left[(1-f_R)I
    + f_R\!\int_{-\infty}^{t}\!R(t-\tau)I(\tau)\,d\tau\right]\mathcal{E} \\[0.3em]
 &- \frac{1}{2}\,\hat T\!\left[\frac{U_i W_{\text{PI}}}{I}\right]\mathcal{E}
   \;-\;\frac{\sigma}{2}\big(1+i\omega_0\tau_c\big)\rho\,\mathcal{E}.
\end{split}
\end{empheq}

\paragraph{Structure.} Terms carrying an explicit $i$ are conservative and act as
phases: diffraction, dispersion, Kerr, and plasma defocusing (the $i\omega_0\tau_c$
part). Terms without $i$ are dissipative: photoionization loss and inverse
bremsstrahlung. This is worth stating because the compact notation
$\partial_z\mathcal{E}=i\hat N\mathcal{E}$ found in the literature is not
literally consistent term by term: with a global $i$ pulled out, the Kerr
contribution inside $\hat N$ is real while the loss terms must themselves carry
an $i$, so that $i\hat N$ reproduces real losses. The boxed
form~\eqref{eq:master} avoids that trap by keeping every factor explicit.

\paragraph{Agreement with the Gaussian derivation.} Every physical prefactor came
out the same: the Kerr coefficient $K_0n_2$, the ionization factor $-\tfrac12$,
and the plasma factor $-\sigma/2$. That is expected --- the two unit systems
must agree --- but it is a real check, because the intermediate expressions look
nothing alike. The place where they genuinely differ is $\chi^{(3)}$, which is a
different number in the two systems, and which is why $n_2$ (not $\chi^{(3)}$)
should always be the quantity carried through the code.

% =====================================================================
\appendix
\section{Gaussian \texorpdfstring{$\leftrightarrow$}{<->} SI dictionary}

\begin{center}
\begin{tabular}{lll}
\toprule
Quantity & Gaussian & SI \\
\midrule
Displacement field & $\mathbf{D}=\mathbf{E}+4\pi\mathbf{P}$
                   & $\mathbf{D}=\varepsilon_0\mathbf{E}+\mathbf{P}$ \\
Source prefactor   & $4\pi/c^2$ & $\mu_0 = 1/\varepsilon_0c^2$ \\
Intensity          & $I=\dfrac{n_0c}{8\pi}|\mathcal{E}|^2$
                   & $I=\dfrac{\varepsilon_0 n_0 c}{2}|\mathcal{E}|^2$ \\[0.6em]
Nonlinear index    & $n_2=\dfrac{12\pi^2\chi^{(3)}}{n_0^2c}$
                   & $n_2=\dfrac{3\chi^{(3)}}{4\varepsilon_0n_0^2c}$ \\[0.6em]
Cross section      & $\sigma=\dfrac{4\pi e^2}{k_0c^2m_{\text{eff}}}
                              \dfrac{\omega_0\tau_c}{1+\omega_0^2\tau_c^2}$
                   & $\sigma=\dfrac{k_0e^2}{n_0^2\omega_0^2\varepsilon_0m_{\text{eff}}}
                              \dfrac{\omega_0\tau_c}{1+\omega_0^2\tau_c^2}$ \\[0.6em]
Kerr term          & $iK_0n_2\hat T^2[I\mathcal{E}]$
                   & $iK_0n_2\hat T^2[I\mathcal{E}]$ \quad(identical) \\
Ionization term    & $-\tfrac12\hat T[U_iW_{\text{PI}}/I]\mathcal{E}$
                   & identical \\
Plasma term        & $-\tfrac{\sigma}{2}(1+i\omega_0\tau_c)\rho\mathcal{E}$
                   & identical \\
\bottomrule
\end{tabular}
\end{center}

\section{Points to check in the Gaussian document}

Collected here for convenience; each is also flagged in place above.

\begin{enumerate}[leftmargin=1.6em]
\item \textbf{The section ``The Lorentz model for STE'' is a verbatim copy of the
Drude section.} Everything after the hydrogenic-exciton introduction repeats the
free-electron derivation word for word, including the sentence ``The density of
current created by a plasma of density $\rho$'' and the final boxed equation.
The Lorentz-oscillator derivation itself (bound resonance at $\omega_{tr}$,
giving $\Delta n \propto +\rho_s/(\omega_{tr}^2-\omega^2)$ with a
\emph{positive} sign, opposite to Drude) is missing. This is the one substantive
gap rather than a typo.

\item \textbf{$\mathbf{J}_{\text{STE}}$ appears in the wave equation but not in
Amp\`ere's law}, where only $\mathbf{J}_{\text{ions}}$ and
$\mathbf{J}_{\text{plasma}}$ are listed.

\item \textbf{$\nabla\times\tilde{\mathbf H}=\nabla\times\tilde{\mathbf B}$}
should be flagged as assuming $\mu=1$.

\item \textbf{Raman drive written twice}, once as $\Omega^2|\mathcal{E}|^2$ and
once as $\Omega^2 I$; the $Q$ substitution is also duplicated.

\item \textbf{Poynting theorem} written $\partial_t u+\nabla S$; should be
$\nabla\!\cdot\!\mathbf{S}$.

\item \textbf{Section heading ``Comparison with''} is incomplete.
\end{enumerate}

\begin{thebibliography}{9}
\bibitem{rothenberg1992}
J.E. Rothenberg,
\textit{Space--time focusing: breakdown of the slowly varying envelope
approximation in the self-focusing of femtosecond pulses},
Opt. Lett. 17, 1340--1342 (1992).

\bibitem{gaeta2000}
A.L. Gaeta,
\textit{Catastrophic collapse of ultrashort pulses},
Phys. Rev. Lett. 84, 3582 (2000).

\bibitem{couairon2007}
A. Couairon, A. Mysyrowicz,
\textit{Femtosecond filamentation in transparent media},
Phys. Rep. 441, 47--189 (2007).

\bibitem{marburger1975}
J.H. Marburger,
\textit{Self-focusing: Theory},
Prog. Quantum Electron. 4, 35--110 (1975).
\end{thebibliography}

\end{document}
